\section*{Chapter \ref{ch:ml:deep-hedging-and-machine-learning-in-pricing} --- Deep Hedging and Machine Learning in Pricing}

\begin{solution}{exo:ml:deep-hedging-and-machine-learning-in-pricing:1}
$\rho(X + c) = \lambda^{-1}\log\E[e^{-\lambda X}e^{-\lambda c}] = \rho(X) - c$. For a constant, $\rho(c) = -c$: holding cash $c$
needs $-c$ more cash to be acceptable, that is, it frees $c$.
\end{solution}

\begin{solution}{exo:ml:deep-hedging-and-machine-learning-in-pricing:2}
$(1.5\times0.001\times100\times0.0998^2)^{1/3} = (0.001494)^{1/3} = 0.114$: a band $\pm0.114$ in delta, 0.229 wide, against the
learned band's 0.055 at the same cost.
\end{solution}

\begin{solution}{exo:ml:deep-hedging-and-machine-learning-in-pricing:3}
Daily hedging with one instrument cannot remove the risk of discrete rebalancing and of stochastic volatility, and the seller
charges for the remaining risk with $\lambda = 1$: 2.454 against 2.254, a risk premium of 0.20.
\end{solution}

\begin{solution}{exo:ml:deep-hedging-and-machine-learning-in-pricing:4}
Under Heston with negative correlation, a fall in the spot comes with a rise in volatility, which raises the call's value; the
variance-minimising hedge ratio is the delta plus a correction for this co-movement (a smaller hedge for a call when $\rho <
0$). The Black--Scholes delta at a fixed implied volatility ignores it; the network, trained on Heston paths, learns it.
\end{solution}

\begin{solution}{exo:ml:deep-hedging-and-machine-learning-in-pricing:5}
Their effects on prices are nearly interchangeable over three maturities: faster mean reversion towards a long-run level and
slower reversion towards a different level can produce similar term structures. Longer maturities, more of them, and
variance-sensitive instruments (variance swaps) separate them.
\end{solution}

\begin{solution}{exo:ml:deep-hedging-and-machine-learning-in-pricing:6}
Stress scenarios lie outside the region the \omterm{def:ml:why-financial-machine-learning-is-different:sets}{test set} sampled, where the surrogate is extrapolating; its test error says nothing
there. Validate on the scenarios themselves, or train on a box that contains them, and check Greeks as well as prices.
\end{solution}

\begin{solution}{exo:ml:deep-hedging-and-machine-learning-in-pricing:7}
The pathwise derivative of $\mathbf1\{S_T > K\}$ with respect to $S_0$ is zero almost surely: the label says the delta is zero
while the true delta is a density at the strike. Pathwise labels need a smooth payoff; the standard remedy is to smooth the
discontinuity (replace the digital by a tight call spread) before differentiating, or to use likelihood-ratio labels.
\end{solution}

\begin{solution}{exo:ml:deep-hedging-and-machine-learning-in-pricing:8}
In a complete market the claim is replicated exactly: a strategy $H$ has gains $G_H = Z - p_{BS}$ path by path. Any strategy
for the seller can be written $H + \delta'$, so $-Z + G_{H + \delta'} = -p_{BS} + G_{\delta'}$ and, by cash invariance,
$\rho(-Z + G_{H+\delta'}) = p_{BS} + \rho(G_{\delta'})$. Taking the infimum over $\delta'$ and subtracting the no-claim term
$\inf_{\delta'}\rho(G_{\delta'})$ leaves $p = p_{BS}$, whatever $\lambda$.
\end{solution}

\begin{solution}{pb:ml:deep-hedging-and-machine-learning-in-pricing:1}
\textbf{Part I.}
\begin{enumerate}
\item A one-month at-the-money call on 100 under Heston, hedged daily with the underlying; the network sees the time to
maturity, the log-moneyness and, through the band, its current holding.
\item 0.123 on average, more than one volatility point of vega (0.115).
\item Not trading is then the default, the band is readable, and training is easier than learning to reproduce the previous
holding exactly.
\item The entropic risk of the P\&L before the premium; $\lambda = 1$ is an absolute risk aversion per unit of price, so a
loss of 1 is weighed $e$ times a loss of 0.
\end{enumerate}
\textbf{Part II.}
\begin{enumerate}[resume]
\item At 0, 5, 10, 20 basis points: deep 2.454, 2.558, 2.658, 2.840; Black--Scholes delta 2.475, 2.608, 2.743, 3.018;
Whalley--Wilmott 2.475, 2.606, 2.712, 2.910.
\item At the money half-way to maturity the learned band is 0.031 wide at 5 basis points and 0.104 at 20, against 0.181 and
0.288; it is narrower at the money and above it, and wider below the money.
\item Lower expected shortfall (3.763 against 3.817) and higher entropic risk (2.671 against 2.658): each objective wins by
its own measure.
\item Its wide band trades least (mean cost 2.424), but the position strays further from the hedge, and the P\&L's standard
deviation (0.703) and tails are the largest of the hedges.
\end{enumerate}
\textbf{Part III.}
\begin{enumerate}[resume]
\item Accuracy per sample: with 1\,024 samples, price error 0.218 against 0.436 and delta error 0.021 against 0.059; at 4\,096,
equal prices and still better deltas.
\item Parameter errors of 1.7\% (initial variance) to 18.7\% (mean-reversion speed) of each range, 13.2 basis points of spot in
re-priced surfaces, 0.46 implied-volatility points on five surfaces.
\item The network's answer as the optimiser's start (and as a check on its result).
\item Outside the training box, on arbitrage constraints, and on market data that no model fits.
\end{enumerate}
\textbf{Part IV.}
\begin{enumerate}[resume]
\item \emph{Hedging with frictions.} The \omterm{def:ml:deep-hedging-and-machine-learning-in-pricing:risk}{indifference price} rises from 2.454 without costs to 2.558, 2.658 and 2.840 at 5, 10
and 20 basis points; the learned band at the money is 0.031 and 0.104 wide at 5 and 20 basis points, against Whalley--Wilmott's
0.181 and 0.288.
\item As a bound on the price at which the desk is compensated for the risk it keeps, not as a quote: the market price and
the book's other positions matter.
\item Out-of-sample paths, other models and parameters, stress paths, and a comparison with the asymptotic band and the delta
on the book's history.
\item The network would also choose how much of the second instrument to hold, hedging volatility risk; the \omterm{def:ml:deep-hedging-and-machine-learning-in-pricing:risk}{indifference
price} would fall towards the Heston value.
\item Retrain or check it against the exact pricer after each recalibration, on a grid covering the new parameters.
\item The hedge is optimal for one volatility dynamics; a different one (higher volatility of variance, a jump) can make it
worse than the delta.
\item It computes the network's own derivatives in differential training, the hedge's gradients in \omterm{def:ml:deep-hedging-and-machine-learning-in-pricing:dh}{deep hedging}, and exact
Greeks of the reference pricer.
\item When trading is worth its cost.
\end{enumerate}
\end{solution}

\begin{solution}{iq:ml:deep-hedging-and-machine-learning-in-pricing:1}
Each rebalancing costs; small deviations from delta add little risk while trading them away costs a spread every time. The
optimal \omterm{def:ml:reinforcement-learning-foundations:mdp}{policy} leaves the position alone while it is close enough to delta and trades back to the band's edge otherwise; the
band widens with costs and narrows with gamma and risk aversion.
\iqlookfor{the cost-risk trade-off and how the band depends on cost, gamma and risk aversion.}
\end{solution}

\begin{solution}{iq:ml:deep-hedging-and-machine-learning-in-pricing:2}
A network \omterm{def:ml:reinforcement-learning-foundations:mdp}{policy} trained on simulated paths: inputs are observable state (time, prices, holdings, other features), outputs
are holdings in the hedging instruments, the loss is a \omterm{def:ml:deep-hedging-and-machine-learning-in-pricing:risk}{convex risk measure} of the terminal hedged P\&L after costs.
\iqlookfor{state, action, loss, and simulation as the training data.}
\end{solution}

\begin{solution}{iq:ml:deep-hedging-and-machine-learning-in-pricing:3}
Train on labels and their derivatives with respect to the inputs, matching the network's derivatives by automatic
differentiation; pathwise derivative labels from Monte Carlo are nearly free and far less noisy relative to their signal, so
fewer samples pin down the function and its slope.
\iqlookfor{the loss and why derivative labels reduce the sample need.}
\end{solution}

\begin{solution}{iq:ml:deep-hedging-and-machine-learning-in-pricing:4}
The premium that leaves the seller indifferent, under a utility or risk measure, between selling with the best hedge and not
selling. With a cash-invariant risk measure it is the risk of the best hedged position without the premium.
\iqlookfor{the definition and its link to cash invariance.}
\end{solution}

\begin{solution}{iq:ml:deep-hedging-and-machine-learning-in-pricing:5}
It is approximate, blind outside its training box, trained on model data that fit perfectly while market data do not, and
gives no fit diagnostics; use it to start the optimiser and to check its result.
\iqlookfor{approximation, extrapolation and model-data mismatch.}
\end{solution}

\begin{solution}{iq:ml:deep-hedging-and-machine-learning-in-pricing:6}
The surrogate is extrapolating or under-trained where prices are tiny and its fit is judged by absolute error; nothing forces
convexity. Train on log prices or implied volatilities, use differential labels including second derivatives, add convexity
constraints or penalties, and sample the wings more densely.
\iqlookfor{the cause and at least two remedies.}
\end{solution}
